%MB9T5
\setcounter{paragraph}{0}
\setcounter{ex}{0}
\PhanI
\begin{ex}%[3P5Y3-1][Loai4]
Thực hiện thí nghiệm Young về giao thoa với ánh sáng đơn sắc màu đỏ ta quan sát được hệ vân giao thoa trên màn. Nếu thay ánh sáng đơn sắc màu đỏ bằng ánh sáng đơn sắc màu lục và các điều kiện khác của thí nghiệm được giữ nguyên thì
\choice
{khoảng vân tăng lên}
{vị trí vân trung tâm thay đổi}
{khoảng vân không thay đổi}
{\True khoảng vân giảm xuống}
\loigiai{Ta có: $i=\dfrac{\lambda D}{a} $ mà ${{\lambda}_{d}}>{{\lambda}_{l}} $ nên khoảng vân giảm xuống.
}
\end{ex} \haidong{20}
\begin{ex}%[3P5Y3-1][Loai4]
Thực hiện thí nghiệm Young về giao thoa với ánh sáng đơn sắc màu vàng ta quan sát được hệ vân giao thoa trên màn. Nếu thay ánh sáng đơn sắc màu vàng bằng ánh sáng đơn sắc màu lam và các điều kiện khác của thí nghiệm được giữ nguyên thì
\choice
{khoảng vân tăng lên}
{\True khoảng vân giảm xuống}
{vị trí vân trung tâm thay đổi}
{khoảng vân không thay đổi}
\loigiai{
Ta có: $i=\dfrac{\lambda D}{a} $ mà ${{\lambda}_{v}}>{{\lambda}_{l}}\Rightarrow {i_v}>{i_l} \Rightarrow$ khoảng vân giảm xuống.
}
\end{ex} \haidong{20}
\loai{Thay đổi môi trường}
\begin{ex}%[3P5B3-1][Loai5]
Trong thí nghiệm Young về giao thoa với ánh sáng đơn sắc. Nếu thực hiện thí nghiệm trong nước thì
\choice
{khoảng vân không đổi}
{tần số thay đổi}
{\True vị trí vân sáng trung tâm không đổi}
{bước sóng không đổi}
\loigiai{
Nếu thực hiện thí nghiệm trên trong nước thì tần số ánh sáng không đổi, chỉ có bước sóng thay đổi
$\Rightarrow $ Khoảng vân cũng thay đổi.
}
\end{ex} \haidong{20}
\loai{Thay đổi bước sóng. Tính khoảng vân}
\begin{ex}%[3P5K3-5][Loai1]
Trong thí nghiệm Young về giao thoa với nguồn sáng đơn sắc, hệ vân trên màn có khoảng vân i. Nếu khoảng cách giữa hai khe còn một nửa và khoảng cách từ hai khe đến màn gấp đôi so với ban đầu thì khoảng vân giao thoa trên màn
\choice
{giảm đi bốn lần}
{không đổi}
{tăng lên hai lần}
{\True tăng lên bốn lần}
\loigiai{
Ta có: $i=\dfrac{\lambda D}{a} $ ta lại có: $i'=\dfrac{\lambda D'}{a'}=\dfrac{\lambda . 2D}{\dfrac{a}{2}}=4i $, vậy khoảng vân tăng 4 lần.
}

\end{ex} \haidong{20}
 \begin{ex}
 	Trong một thí nghiệm Young về giao thoa ánh sáng với ánh sáng đơn sắc có bước sóng $\lambda_{1}=540\ nm$ thì thu được hệ vân giao thoa trên màn quan sát có khoảng vân $i_{1}=0,36\ mm$. Khi thay ánh sáng trên bằng ánh sáng đơn sắc có bước sóng $\lambda_{2}=600\ nm$ thì thu được hệ vân giao thoa trên màn quan sát có khoảng vân bằng
 	\choice
 	{\True $0,4\ mm$}
 	{$0,5\ mm$}
 	{$0,6\ mm$}
 	{$0,7\ mm$}
 	\loigiai{  }
 \end{ex}
\loai{Thay đổi bước sóng. Tính bước sóng}
\begin{ex}%[3P5K3-5][Loai2]
Trong thí nghiệm Young: ánh sáng được dùng là ánh sáng đơn sắc có bước sóng $\lambda =0,52\ \mu m$. Thay ánh sáng trên bằng ánh sáng đơn sắc có bước sóng $\lambda '$ thì khoảng vân tăng 1,3 lần. Bước sóng $\lambda'$ có giá trị 
\choice
{4 $\ \mu m$}
{0,4$\ \mu m$}
{6,8$\ \mu m$}
{\True 0,68$\ \mu m$}
\loigiai{
\begin{itemize}
\item Ta có: $\dfrac{i}{i'}=\dfrac{\lambda}{\lambda '}\Rightarrow \lambda '= 0,68\ \mu m$.
\end{itemize}
}
\end{ex} \haidong{20}
\begin{ex}%[3P5K3-5][Loai1]
Trong thí nghiệm Young về giao thoa ánh sáng, khi dùng ánh sáng đơn sắc có bước sóng 0,4\microm\ thì khoảng vân đo được là 0,2\ mm. Nếu dùng ánh sáng đơn sắc có bước sóng 0,6\microm mà vẫn giữ nguyên khoảng cách giữa hai khe và từ hai khe tới màn thì khoảng vân là
\choice
{0,4\ mm}
{0,2\ mm}
{0,6\ mm}
{\True 0,3\ mm}
\loigiai{
\begin{itemize}
\item Từ công thức tính khoảng vân: \bth{i=\dfrac{\lambda.D}{a}}.
\begin{itemize}
\item Vì D không đổi nên i tỉ lệ thuận với \bth{\lambda}.
\item \bth{\lambda} tăng từ 0,4\microm\ lên 0,6\microm, tức là tăng 1,5 lần.
\end{itemize}
\item Suy ra khoảng vân cũng tăng 1,5 lần: \bth{i'=1,5i=0,2.1,5=0,3\ mm}.
\end{itemize}
}
\end{ex} \haidong{20}

\begin{ex}%CD
Người ta đặt lần lượt các tấm kính lọc trước nguồn phát ánh sáng trắng trong thí nghiệm giao thoa Young. Lúc đầu, khi dùng kính lọc màu đỏ có bước sóng 600 nm thì khoảng vân đo được là 2,40 mm. Khi dùng kính lọc màu lam thì khoảng vân đo được là 1,80 mm. Bước sóng của ánh sáng đi qua kính lọc màu lam là
\choice
{\True $450\ nm$}
{$550\ nm$}
{$650\ nm$}
{$750\ nm$}
\loigiai{
\begin{itemize}
\item $\dfrac{i_{\ct{đỏ}}}{i_{\ct{lam}}}=\dfrac{\lambda_{\ct{đỏ}}}{\lambda_{\ct{lam}}}\Rightarrow \lambda_{lam}=450\ nm$.
\end{itemize}
}
\end{ex} \haidong{20}
\begin{ex}%[3P5K3-5][Loai2]
Trong một thí nghiệm Young về giao thoa ánh sáng, khoảng cách giữa hai khe là 2 mm, khoảng cách từ hai khe đến màn là 1 m. Khi dùng ánh sáng đơn sắc có bước sóng $\lambda_1 $ để làm thí nghiệm thì đo được khoảng cách giữa 5 vân sáng liên tiếp nhau là 0,8 mm. Thay bức xạ có bước sóng $\lambda_1 $ bằng bức xạ có bước sóng $\lambda_2>\lambda_1 $ thì tại vị trí của vân sáng bậc 3 của bức xạ bước sóng $\lambda_1 $ ta quan sát được một vân sáng của bức xạ có bước sóng $\lambda_2 $. Giá trị $\lambda_2 $ là?
\choice
{ $0,38\ \mu m $}
{ $0,40\ \mu m $}
{ $0,53\ \mu m $}
{\True $0,6\ \mu m $}
\loigiai{
\begin{itemize}
\item Ta có: $4{i_1}=0,8\Rightarrow {i_1}=\dfrac{\lambda_1D}{a}=0,2\Rightarrow \lambda_1=0,4\ \mu m$.
\item Tại vị trí của vân sáng bậc 3 của bức xạ bước sóng $\lambda_1 $ ta quan sát được một vân sáng của bức xạ có bước sóng $\lambda_2 \Rightarrow 3\lambda_1=k\lambda_2 $ mà $\lambda_2>\lambda_1 $ $\Rightarrow k<3\Rightarrow k=1;2$.
\item Với $k = 1 \Rightarrow \lambda_2=1,2\ \mu m $ (loại vì ngoài khoảng nhìn thấy).
\item Với $k = 2 \Rightarrow \lambda_2=0,6\ \mu m $ (thỏa mãn).
\end{itemize}
}
\end{ex} \haidong{20}

\loai{Thay đổi bước sóng. Tính tổng hợp}
\begin{ex}%[3P5B3-5][Loai3]
Trong thí nghiệm Young, hai khe được chiếu ánh sáng có bước sóng ${{\lambda}_{1}}=500\ nm $. Xét một điểm trên màn mà hiệu đường đi tới hai nguồn sáng là 0,75 $\mu m $, tại điểm này quan sát được gì nếu thay ánh sáng trên bằng ánh sáng có bước sóng ${{\lambda}_{2}}=750\ nm $
\choice
{\True từ vân tối thành vân sáng}
{từ vân sáng thành vân tối}
{đều cho vân sáng}
{đều cho vân tối}
\loigiai{
\begin{itemize}
\item Khi được chiếu bước sóng ${{\lambda}_{1}} $ thì ${k_1}=\dfrac{{d_2}-{d_1}}{{{\lambda}_{1}}}=\dfrac{0,75}{0,5}=1. 5 $, ${k_1} $ bán nguyên nên lúc này tại điểm xét là vân tối. 
\item Khi thay bằng ${{\lambda}_{2}} $ thì ${k_2}=\dfrac{0,75}{0,75}=1 $ nguyên nên điểm xét là vân sáng.
\end{itemize}
}
%<MyLT>
\end{ex} \haidong{20}
\loai{Thay đổi a}
\begin{ex}%[3P5K3-5][Loai4]
Trong thí nghiệm Young về giao thoa với ánh sáng đơn sắc xác định, thì tại điểm M trên màn quan sát là vân sáng bậc 5. Sau đó giảm khoảng cách giữa hai khe một đoạn bằng 0,2 mm thì tại M trở thành vân tối thứ 5 so với vân sáng trung tâm. Ban đầu khoảng cách giữa hai khe là
\choice
{2,2 mm}
{1,2 mm}
{\True 2 mm}
{1 mm}
\loigiai{
\begin{itemize}
\item Ta có: $x_M=5\dfrac{\lambda D}{a}=4,5\dfrac{\lambda D}{a-0,2}\Rightarrow \dfrac{5}{a}=\dfrac{4,5}{a-0,2}\Rightarrow a=2\ mm. $
\end{itemize}
}

\end{ex} \haidong{20}
\begin{ex}
Trong thí nghiệm Young về giao thoa với ánh sáng đơn sắc xác định. Ban đầu khoảng cách giữa hai khe là $a$ thì tại điểm $M$ trên màn quan sát có vân sáng bậc 7. Sau đó, người ta giảm khoảng cách giữa hai khe đi một đoạn bằng $0,2\ mm$ thì tại $M$ trở thành vân tối thứ 6. Khoảng cách ban đầu giữa hai khe là 
\choice
{$0,33\ mm$}
{$0,63\ mm$}
{\True $0,93\ mm$}
{$0,13\ mm$}
\loigiai{ }
\end{ex}
\begin{ex}%[3P5G3-5][Loai4]
Trong thí nghiệm Young, nguồn S phát bức xạ đơn sắc $\lambda$, màn quan sát cách mặt phẳng hai khe một khoảng không đổi D, khoảng cách giữa hai khe $S_1S_2=a$ có thể thay đổi (nhưng $S_1$ và $S_2$ luôn cách đều S). Xét điểm M trên màn, lúc đầu là vân sáng bậc 4, nếu lần lượt giảm hoặc tăng khoảng cách $S_1S_2$ một lượng $\Delta a$ thì tại đó là vân sáng bậc k và bậc 3k. Nếu tăng khoảng cách $S_1S_2$ thêm $2\Delta a$ thì tại M là
\choice
{vân tối thứ 9}
{vân sáng bậc 9}
{vân sáng bậc 7}
{\True vân sáng bậc 8}
\loigiai{\begin{itemize}
\item Khi khoảng cách 2 khe tới màn là a thì tại M là vân sáng bậc 4 nên $x_M=4.\dfrac{\lambda D}{a}\quad (1)$.
\item Nếu lần lượt giảm hoặc tăng khoảng cách $S_1S_2$ một lượng $\Delta a$ thì tại đó là vân sáng bậc k và bậc 3k nên
$\begin{cases}
x_M=k.\dfrac{\lambda D}{a-\Delta a} \\
x_M=3k.\dfrac{\lambda D}{a+\Delta a} \\
\end{cases}\Rightarrow \dfrac{k}{a-\Delta a}=\dfrac{3k}{a+\Delta a}\Rightarrow a=2.\Delta a$.
\item Nếu tăng khoảng cách $S_1S_2$ thêm $2\Delta a$ thì tại M là:
$x_M=k'.\dfrac{\lambda D}{a+2\Delta a}=k'.\dfrac{\lambda D}{a+a}=\dfrac{1}{2}{k}'.\dfrac{\lambda D}{a}\quad (2)$.
\item So sánh với $(1)$ ta có: $x_M=\dfrac{1}{2}{k}'.\dfrac{\lambda D}{a}=4.\dfrac{\lambda D}{a}\Rightarrow k'=8\Rightarrow $ Tại M khi đó là vân sáng bậc 8.
\end{itemize}
}
%<MyLT>
\end{ex} \haidong{20}
\loai{Dịch màn ra xa. Tìm bước sóng}

\begin{ex}%[3P5K3-5][Loai5]
Trong thí nghiệm giao thoa Young, khoảng cách hai khe là 1\ mm. Giao thoa thực hiện với ánh sáng đơn sắc có bước sóng $\lambda$ thì tại điểm M có tọa độ 1,2\ mm là vị trí vân sáng bậc 4. Nếu dịch màn xa thêm một đoạn 25 cm theo phương vuông góc với mặt phẳng hai khe thì tại M là vị trí vân sáng bậc 3. Bước sóng có giá trị bằng
\choice
{\True $0,4\ \mu m$}
{$0,48\ \mu m$}
{$0,45\ \mu m$}
{$0,44\ \mu m$}
\loigiai{\begin{itemize}
\item $\left\{\begin{aligned}& x_M=4\dfrac{\lambda D}{a}\Rightarrow \dfrac{\lambda D}{a}=\dfrac{x_M}{4} \\
& x_M=3\dfrac{\lambda (D+0,25)}{a}=3\dfrac{\lambda D}{a}+0,75.\dfrac{\lambda}{a}
\end{aligned}\right.\Rightarrow \lambda =0,4.10^{-6}\ m $.
\end{itemize}
}
\end{ex} \haidong{20}
\begin{ex}%[3P5K3-5][Loai5]
Thí nghiệm giao thoa Young với ánh sáng đơn sắc có bước sóng $\lambda $, khoảng cách giữa hai khe 0,5 mm. Ban đầu, tại M cách vân trung tâm 1 mm người ta quan sát được vân sáng bậc 2. Giữ cố định màn chứa hai khe, di chuyển từ từ màn quan sát ra xa và dọc theo đường thẳng vuông góc với mặt phẳng chứa hai khe một đoạn $\dfrac{50}{3} $ cm thì thấy tại M chuyển thành vân tối thứ 2 kể từ vân trung tâm. Bước sóng $\lambda $ có giá trị là
\choice
{0,60 $\mu m $}
{\True 0,50 $\mu m $}
{0,40 $\mu m $}
{0,64 $\mu m $}
\loigiai{\begin{itemize}
\item Ta có: $OM=2i=1,5i'=1\Rightarrow i=0,5\ mm;\ i'=\dfrac{2}{3}\ mm $.
\item Lại có: $\dfrac{i}{i'}=\dfrac{D}{D'}\Rightarrow \dfrac{3}{4}=\dfrac{D}{D+50/3}\Rightarrow D=50\ cm\ =0,5\ m $ $\Rightarrow \lambda =\dfrac{ia}{D}=\dfrac{0,5. 0,5}{0,5}=0,5\ \mu m $.
\end{itemize}
}
%<MyLT>
\end{ex} \haidong{20}
\loai{Dịch màn ra xa. Tính tổng hợp}
\begin{ex}%[3P5K3-5][Loai6]
Trong thí nghiệm Young về giao thoa ánh sáng đơn sắc với hai khe sáng cách màn quan sát 1,375 m thì tại điểm M trên màn quan sát được vân sáng bậc 5. Để quan sát được vân tối thứ 6 tại điểm M nói trên thì phải tịnh tiến màn theo phương vuông góc với nó một đoạn
\choice
{\True 0,125 m}
{0,25 m}
{0,2 m}
{0,115 m}
\loigiai{
\begin{itemize}
\item Để tại M là vân tối thứ 6 thì $i'=\dfrac{OM}{5,5}=\dfrac{5i}{5,5}\Rightarrow \dfrac{i'}{i}=\dfrac{5}{5,5}=\dfrac{D'}{D}\Rightarrow \dfrac{D-x}{D}=\dfrac{5}{5,5}\Rightarrow x=0,125\ m. $
\end{itemize}
}
\end{ex} \haidong{20}

\begin{ex}%[3P5K3-5][Loai6]
Trong thí nghiệm Young về giao thoa ánh sáng, hai khe được chiếu bằng ánh sáng đơn sắc. Từ vị trí ban đầu, nếu tịnh tiến màn quan sát một đoạn 50 cm lại gần mặt phẳng chứa hai khe thì khoảng vân mới thay đổi một lượng bằng 250 lần bước sóng. Tính khoảng cách giữa hai khe hẹp.
\choice
{20 mm}
{\True 2 mm}
{1 mm}
{3 mm}
\loigiai{
\begin{itemize}
\item Theo bài ra $i=i'+250\lambda \Rightarrow \dfrac{\lambda D}{a}=\dfrac{\lambda D'}{a}+250\lambda \Rightarrow \dfrac{D}{a}=\dfrac{D'}{a}+250\Rightarrow a=\dfrac{D-D'}{250}=\dfrac{50}{250}=0,2\ cm $.
\end{itemize}
}

\end{ex} \haidong{20}
\loai{Dịch chuyển màn lại gần}
\begin{ex}%KNTT
Trong thí nghiệm Young về giao thoa ánh sáng, hai khe được chiếu bằng ánh sáng đơn sắc. Khoảng cách giữa hai khe bằng 0,6 mm. Khoảng vân trên màn quan sát đo được là 1 mm. Từ vị trí ban đầu, nếu tịnh tiến màn quan sát một đoạn 25 cm lại gần mặt phẳng chứa hai khe thì khoảng vân mới trên màn là 0,8 mm. Bước sóng của ánh sáng dùng trong thí nghiệm là
\choice
{$0,56\ \mu m$}
{\True $0,48\ \mu m$}
{$0,66\ \mu m$}
{$0,74\ \mu m$}
\loigiai{
\begin{itemize}
\item Khoảng cách giữa hai khe là 0,6 mm và khoảng vân trên màn quan sát được là 1 mm.
\item  Ta có tỉ lệ: $\dfrac{i_1}{i_2}=\dfrac{D_1}{D_2}=\dfrac{D}{D-0,25}\Rightarrow \dfrac{D}{D-0,25}=\dfrac{1}{0,8}\Rightarrow D=1,25\ m$
\item Vậy bước sóng của ánh sáng dùng trong thí nghiệm là: $\lambda=\dfrac{ia}{D}=0,48\ \mu m$.
\end{itemize}
}
\end{ex} \haidong{20}
\setcounter{ex}{0}
\PhanII
\begin{ex}
Một nhóm học sinh thực hành đo bước sóng ánh sáng đơn sắc bằng thí nghiệm giao thoa Young. Khoảng cách giữa hai khe $S_1$ và $S_2$ là $a=0,10\ mm$. Đo khoảng cách từ màn tới hai khe Young là $D=700\ mm$. Thí nghiệm cho thấy trên màn quan sát có $7$ vân sáng. Dùng thước kẹp, học sinh đo được khoảng cách trung bình giữa hai tâm của hai vân sáng ngoài cùng là $L=30,12\ mm$.
\choiceTF[t]
{\True Bước sóng trung bình đo được xấp xỉ $717\ nm$}
{Nếu giữ nguyên khoảng cách giữa hai khe và dịch màn theo phương vuông góc với mặt phẳng hai khe, ra xa thêm $0,3\ m$ thì trên màn chỉ quan sát thấy 3 vân sáng liên tiếp}
{\True Khoảng vân xấp xỉ $5,02\ mm$}
{Trong khoảng 7 vân sáng rõ liên tiếp trên màn có 7 khoảng vân}
\loigiai{
\begin{itemchoice}
\itemch 
\itemch 
\itemch 
\itemch 
\end{itemchoice}
}
\end{ex}
\begin{ex}
Chiếu bức xạ đơn sắc bước sóng $0,5~ \mu m$ vào hai khe Young cách nhau 1 mm. Màn quan sát đặt cách hai khe Young 2 m.
\choiceTF[t]
{Vân trung tâm là vân sáng}
{\True Khoảng vân trên màn là $i = 1\ mm$}
{\True Khoảng cách từ vân sáng bậc 3 đến vân tối thứ 1 ở hai bên vân trung tâm là $3,5 \  mm$}
{\True Tăng khoảng cách từ hai khe đến màn thêm $50\ cm$ thì khoảng vân tăng một lượng $0,25 \  mm$}
\loigiai{
\begin{itemchoice}
\itemch Vân trung tâm là vân sáng bậc k=0.
\itemch $i = \dfrac{\lambda D}{a} = \dfrac{0.5 . 10^{-6} . 2}{1 . 10^{-3}} = 10^{-3} \  m = 1 \  mm$.
\itemch Vân sáng bậc 3 cách vân trung tâm $3i = 3 \  mm$. Vân tối thứ 1 cách vân trung tâm $0.5i = 0.5 \  mm$. Khoảng cách là $3.5 \  mm$.
\itemch $i' = \dfrac{\lambda (D + 0.5)}{a} = \dfrac{0.5 . 10^{-6} . (2 + 0.5)}{1 . 10^{-3}} = \dfrac{0.5 . 2.5 . 10^{-3}}{1} = 1.25 \  mm$. $\Delta i = i' - i = 1.25 - 1 = 0.25 \  mm$.
\end{itemchoice}
}
\end{ex}
\begin{ex}
Trong thí nghiệm Young sử dụng ánh sáng đơn sắc, hai khe cách nhau $0,2\ mm$, màn đặt cách hai khe $1\ m$. Khi dùng một nguồn sáng đơn sắc, đo được khoảng cách giữa hai vân sáng liên tiếp là $3\ mm$.
\choiceTF[t]
{\True Bước sóng ánh sáng đơn sắc sử dụng trong thí nghiệm là $600\ nm$}
{Khoảng cách từ vân sáng bậc hai đến vân tối thứ tư ở cùng bên là $3\ mm$}
{Khi sử dụng đồng thời với một ánh sáng có bước sóng $\lambda' = 450\ nm$, vị trí vân sáng của 2 hệ vân trùng nhau nằm gần vân sáng trung tâm nhất cách vân trung tâm $13,5\ mm$}
{\True Trong cùng điều kiện thí nghiệm, khoảng vân $i$ của ánh sáng đơn sắc có bước sóng $450\ nm$ là $2,25\ mm$}
\loigiai{
\begin{itemchoice}
\itemch $i = 3\ mm$, $a = 0,2\ mm$, $D = 1\ m \Rightarrow \lambda = \dfrac{a i}{D} = \dfrac{0,2 . 10^{-3} . 3 . 10^{-3}}{1} = 600\ nm$.
\itemch $x_{s2} = 2i = 6\ mm$, $x_{t4} = 3,5i = 10,5\ mm \Rightarrow \Delta x = 4,5\ mm$.
\itemch Theo tính toán, vị trí trùng vân sáng gần nhất ứng với $k_1 = 3, k_2 = 4 \Rightarrow x = 3i = 9\ mm$
\itemch Vì $i' = \dfrac{\lambda' D}{a} = \dfrac{450 . 10^{-9} . 1}{0,2 . 10^{-3}} = 2,25\ mm$
\end{itemchoice}
}
\end{ex}
\begin{ex}
Trong thí nghiệm Young về giao thoa ánh sáng, nguồn sáng phát ánh sáng đơn sắc có bước sóng $\lambda=0,4\ \mu m$. Khoảng cách giữa hai khe là $a=0,5\ mm$, khoảng cách từ hai khe đến màn là $D=2\ m$. M và N là hai điểm khác phía với vân trung tâm, cách vân trung tâm lần lượt $6,5\ mm$ và $18\ mm$.
\choiceTF[t]
{Khoảng vân trên màn là $0,8\ mm$}
{Khoảng cách giữa vân sáng thứ 3 và vân tối thứ 5 ở cùng bên so với vân trung tâm bằng $3,2\ mm$}
{\True Trên đoạn MN quan sát được 14 vân tối}
{\True Nếu tăng khoảng cách từ hai khe đến màn thêm $25\ cm$ thì khoảng vân đo được bằng $1,8\ mm$}
\loigiai{
\begin{itemchoice}
\itemch 
\itemch 
\itemch 
\itemch 
\end{itemchoice}
}
\end{ex}


\setcounter{ex}{0}
\PhanIII
\begin{ex}
Trong thí nghiệm Young, khi sử dụng ánh sáng đỏ có bước sóng $0,72\ \mu m$ thì khoảng vân thu được là $1,2\ mm$ và khi sử dụng một ánh sáng khác có bước sóng $\lambda\ (\mu m)$ thì khoảng vân thu được bằng $0,9\ mm$. Giá trị của $\lambda$ là bao nhiêu $\mu m$?
\shortans[oly]{0,54}
\loigiai{ }
\end{ex}
\begin{ex}
Trong thí nghiệm của Young, người ta dùng ánh sáng đơn sắc có bước sóng $\lambda=0,75\ \mu m$. Nếu thay ánh sáng trên bằng ánh sáng đơn sắc có bước sóng $\lambda'$ thì thấy khoảng vân giao thoa giảm đi 1,5 lần. Giá trị của $\lambda'$ bằng bao nhiêu $\mu m$?\\\shortans[oly]{0,5}
\loigiai{
0,5
}
\end{ex} \haidong{20}
\begin{ex}
Thực hiện thí nghiệm Young về giao thoa với ánh sáng đơn sắc có bước sóng $\lambda$. Khoảng cách giữa hai khe hẹp là $1 \  mm$. Trên màn quan sát, tại điểm M cách vân trung tâm $4,2 \  mm$ có vân sáng bậc 5. Giữ cố định các điều kiện khác, di chuyển dần màn quan sát dọc theo đường thẳng vuông góc với mặt phẳng chứa hai khe ra xa cho đến khi vân giao thoa tại M chuyển thành vân tối thứ tư thì khoảng dịch màn là $0,6 \   m$. Bước sóng $\lambda$ của ánh sáng này bằng $x\ \mu m$. Tìm x.
\shortans[oly]{0,6}
\loigiai{
Vị trí vân sáng bậc 5 là 
$$x_M = x_{s5} = 5i = 4,2 \  mm$$
Từ đó, ta tính được khoảng vân ban đầu:
$$i = \dfrac{4,2}{5} = 0,84 \  mm$$
Khi màn dịch chuyển ra xa một khoảng $0,6 \   m$, khoảng vân mới là $i'$, và tại M là vân tối thứ tư, vị trí của nó là 
$$x_M = x_{t4} = 3,5 i'$$
Vì vị trí M không đổi, ta có:
$$5i = 3,5 i'$$
Mà khoảng vân $i = \dfrac{\lambda D}{a}$ và $i' = \dfrac{\lambda (D + 0,6)}{a}$. Thay vào phương trình trên:
$$5 \dfrac{\lambda D}{a} = 3,5 \dfrac{\lambda (D + 0,6)}{a}$$
Rút gọn $\dfrac{\lambda}{a}$, ta được:
$$5D = 3,5(D + 0,6)$$
$$5D = 3,5D + 2,1$$
$$1,5D = 2,1$$
$$D = \dfrac{2,1}{1,5} = 1,4 \   m$$
Bước sóng $\lambda$ được tính bằng công thức:
$$\lambda = \dfrac{ai}{D} = \dfrac{1 . 10^{-3} \   m . 0,84 . 10^{-3} \   m}{1,4 \   m} = 0,6 . 10^{-6} \   m = 0,6 \  \mu m$$
}
\end{ex}
\begin{ex}
Trong thí nghiệm Young về giao thoa ánh sáng,  hai khe được chiếu bằng ánh sáng đơn sắc $\lambda$, màn quan sát cách mặt phẳng hai khe một khoảng không đổi $D$, khoảng cách giữa hai khe có thể thay đổi (nhưng $S_1$ và $S_2$ luôn cách đều $S$). Xét điểm $M$ trên màn, lúc đầu là vân tối thứ $4$, nếu lần lượt giảm hoặc tăng khoảng cách $S_1S_2$ một lượng $\Delta a$ thì tại đó là vân sáng bậc $k$ và bậc $3k$. Nếu tăng khoảng cách $S_1S_2$ thêm $2\Delta a$ thì tại $M$ là vân sáng bậc mấy?
\shortans[oly]{7}
\loigiai{
Tính theo công thức bậc vân khi thay đổi khoảng cách giữa hai khe, suy ra bậc vân mới là $7$.
}
\end{ex}
\begin{ex}
Trong thí nghiệm giao thoa Young, khoảng cách giữa hai khe và khoảng cách từ hai khe đến màn lần lượt là $a$ và $D$. Sử dụng nguồn laser đỏ có bước sóng $\lambda$, trên màn quan sát đo được khoảng cách giữa 3 vạch sáng màu đỏ liên tiếp là $1,2\ cm$. Nếu thay thế bằng ánh sáng có bước sóng $\lambda' = 0,6\lambda$ thì khoảng cách giữa 3 vạch sáng trên màn quan sát bằng bao nhiêu cm (làm tròn kết quả đến chữ số hàng phần mười)?
\shortans[oly]{0,7}
\loigiai{}
\end{ex}
\begin{ex}
Thực hiện thí nghiệm Young về giao thoa với ánh sáng đơn sắc có bước sóng $\lambda$. Khoảng cách giữa hai khe hẹp là $1 \  mm$. Trên màn quan sát, tại điểm M cách vân trung tâm $4,2\  mm$ có vân sáng bậc 5. Giữ cố định các điều kiện khác, di chuyển dần màn quan sát dọc theo đường thẳng vuông góc với mặt phẳng chứa hai khe ra xa cho đến khi vân giao thoa tại M chuyển thành vân tối thứ tư thì khoảng dịch màn là $0,6 \   m$. Bước sóng $\lambda$ của ánh sáng này bằng bao nhiêu $\mu m$?
\shortans[oly]{0,6}
\loigiai{
Vị trí vân sáng bậc 5 là 
$$x_M = x_{s5} = 5i = 4,2 \  mm$$
Từ đó, ta tính được khoảng vân ban đầu:
$$i = \dfrac{4,2}{5} = 0,84 \  mm$$
Khi màn dịch chuyển ra xa một khoảng $0,6 \   m$, khoảng vân mới là $i'$, và tại M là vân tối thứ tư, vị trí của nó là 
$$x_M = x_{t4} = 3,5 i'$$
Vì vị trí M không đổi, ta có:
$$5i = 3,5 i'$$
Mà khoảng vân $i = \dfrac{\lambda D}{a}$ và $i' = \dfrac{\lambda (D + 0,6)}{a}$. Thay vào phương trình trên:
$$5 \dfrac{\lambda D}{a} = 3,5 \dfrac{\lambda (D + 0,6)}{a}$$
Rút gọn $\dfrac{\lambda}{a}$, ta được:
$$5D = 3,5(D + 0,6)$$
$$5D = 3,5D + 2,1$$
$$1,5D = 2,1$$
$$D = \dfrac{2,1}{1,5} = 1,4 \   m$$
Bước sóng $\lambda$ được tính bằng công thức:
$$\lambda = \dfrac{ai}{D} = \dfrac{1 . 10^{-3} \   m . 0,84 . 10^{-3} \   m}{1,4 \   m} = 0,6 . 10^{-6} \   m = 0,6 \  \mu m$$
}
\end{ex}

